\begin{graphillustration}{Пример из текста: многочлен дерева на четырёх вершинах показывает старший коэффициент, число рёбер и число компонент.}{fig:20-coefficients}
\begin{tikzpicture}[font=\small]
\node[gi vertex] (1) at (0,0) {$1$};
\node[gi vertex] (2) at (1,1) {$2$};
\node[gi vertex] (3) at (2,0) {$3$};
\node[gi vertex] (4) at (3,1) {$4$};
\draw[gi edge] (1) -- (2);
\draw[gi edge] (2) -- (3);
\draw[gi edge] (3) -- (4);
\node[gi note] at (1.5,-0.8) {$n=4,\ m=3,\ c=1$};
\node[gi note] at (7.4,0.4) {$\chi_G(k)=\underbrace{1}_{b_0}k^4\underbrace{-3}_{-m}k^3+3k^2\underbrace{-1}_{\ne0}k$};
\node[gi note] at (7.4,-0.9) {Ненулевые степени: $4,3,2,1$; младшая — $c(G)=1$.};
\end{tikzpicture}
\end{graphillustration}
